\section{复杂构型、锁型、伪椭圆与连分式算法}\label{ic:indefinite_algebra_complete}

\begin{examtip}本节将隐藏的微分结构称为锁型，并用复合微分、商导数、线性微分方程和代数恒等式解释。连分式部分给出有限步骤及最后的多项式验算。\end{examtip}

\subsection{二项微分的三条有理化路线}\label{ic:indefinite_algebra_complete:binomial}\begin{method}{二项微分的三条有理化路线}\textbf{识别信号} $\int x^m(a+bx^n)^p\dd x$，指数有理，$n\ne0$；先在$x>0,a>0,a+bx^n>0$的实支讨论。\par \textbf{操作步骤} 若$p$为整数，用共同根次有理化；若$q=\frac{m+1}{n}$为整数，令$u=a+bx^n$后取分数幂公分母；若$p+q$为整数，令$u=b+ax^{-n}$后同样消去分数幂。\par \textbf{疑难点与条件} 这里证明的是三条可执行的充分路线；条件均不满足时，不凭这一张卡就断言非初等。\end{method}

\begin{theorem}{第三条路线的变换公式}令$q=\frac{m+1}{n}$、$k=p+q$。换$u=b+ax^{-n}$后，$x^m(a+bx^n)^p\dd x=-\frac{a^k u^p(u-b)^{-k-1}\dd u}{n}$。若$k\in\mathbb Z$，再令$u=t^d$，$d$为$p$的分母，就得到有理式。\end{theorem}\begin{proof}写$x^n=\frac{a}{u-b}$，于是$x^m\dd x=-(\frac{a^q}{n})(u-b)^{-q-1}\dd u$；又$(a+bx^n)^p=a^pu^p(u-b)^{-p}$，相乘即得。第二条路线类似：令$u=a+bx^n$后，积分是常数倍$(u-a)^{q-1}u^p\dd u$，$q$整数时取公分母即可有理化。\end{proof}

\begin{example}{竞赛风格｜整数$q$的非平方根积分}{难度：$\star\star\star$}\label{ic:ex:indefinite_algebra_complete:binomialq}\textbf{题目} 求$\int x^5(1+x^3)^{\frac{1}{3}}\dd x$，$x>0$。\par \textbf{思路} $q=\frac{5+1}{3}=2$，先取$u=1+x^3$。\par \textbf{详细解答} $x^5\dd x=\frac{(u-1)\dd u}{3}$，所以积分$\frac13\int(u^{\frac{4}{3}}-u^{\frac{1}{3}})\dd u=\frac{u^{\frac{7}{3}}}{7}-\frac{u^{\frac{4}{3}}}{4}+C$。也可直接设$t=(1+x^3)^{\frac{1}{3}}$，积分化为$(t^6-t^3)\dd t$。\par \textbf{求导验算} 求导为$x^2[(1+x^3)^{\frac{4}{3}}-(1+x^3)^{\frac{1}{3}}]=x^5(1+x^3)^{\frac{1}{3}}$。\par \textbf{疑难点与方法比较} 不必看到立方根就引入特殊函数。\end{example}

\begin{example}{竞赛风格｜整数$p+q$路线双解}{难度：$\star\star\star$}\label{ic:ex:indefinite_algebra_complete:binomialsum}\textbf{题目} 求$\int\frac{\sqrt{1+x^2}\dd x}{x^2}$，$x>0$。\par \textbf{思路} $p=\frac{1}{2},q=-\frac{1}{2}$，其和为0。\par \textbf{详细解答} 令$t=\sqrt{1+x^{-2}}=\frac{\sqrt{1+x^2}}{x}>1$，$\dd t=-\frac{\dd x}{x^2\sqrt{1+x^2}}$，原式$=-\int \frac{t^2}{t^2-1}\dd t=-t+\frac12\ln|\frac{t+1}{t-1}|+C$。另一法分部取$u=\sqrt{1+x^2},\dd v=x^{-2}\dd x$，得$-\frac{\sqrt{1+x^2}}{x}+\ln(x+\sqrt{1+x^2})+C$。\par \textbf{求导验算} 第二式求导为$\frac{1}{\sqrt{1+x^2}}+\frac{1}{x^2\sqrt{1+x^2}}=\frac{\sqrt{1+x^2}}{x^2}$。\par \textbf{疑难点与方法比较} 二项路线给系统换元，分部法在低次时更短。\end{example}

\subsection{Hermite约化：模多项式算法}\label{ic:indefinite_algebra_complete:hermitemod}\begin{method}{Hermite约化：模多项式算法}\textbf{识别信号} $\frac{A}{V^kW}$，$k\ge2$，$V$平方自由且与$W$互素。\par \textbf{操作步骤} 求次数小于$\deg V$的$B$，使$-(k-1)WV'B\equiv A\pmod V$；由扩展Euclid求逆。令$C=\frac{A-WVB'+(k-1)WV'B}{V}$，减去$(\frac{B}{V^{k-1}})'$，分母幂数降1。重复直到平方自由。\par \textbf{疑难点与条件} 模逆存在依赖$\gcd(V,WV')=1$；不满足时先因式分解及平方自由分解。\end{method}

\begin{theorem}{Hermite单步约化的证明}按上法定义$B,C$，则$\frac A{V^kW}=\left(\frac B{V^{k-1}}\right)'+\frac C{V^{k-1}W}$。\end{theorem}\begin{proof}右侧通分后分子为$WVB'-(k-1)WV'B+CV=A$。模同余使$C$确为多项式，分母中$V$的次数下降，因而重因子部分有限步终止。最后的平方自由真分式通过实或复根分解产生对数与反正切。\end{proof}

\begin{example}{竞赛风格｜一次完整模法约化}{难度：$\star\star\star$}\label{ic:ex:indefinite_algebra_complete:hermitefull}\textbf{题目} 求$\int\frac{\dd x}{(x^2+1)^2(x+1)}$。\par \textbf{思路} $V=x^2+1,W=x+1,k=2$。\par \textbf{详细解答} 模$V$有$-WV'\equiv2-2x$，其逆为$\frac{1+x}{4}$，故$B=\frac{1+x}{4}$。计算$C=\frac{x+3}{4}$，于是原式导数部分为$[\frac{x+1}{4V}]'$，剩余$\frac{x+3}{4V(x+1)}=\frac{1}{4(x+1)}+\frac{-x+2}{4V}$。积分得$\frac{x+1}{4(x^2+1)}+\frac14\ln|x+1|-\frac18\ln(x^2+1)+\frac12\arctan x+C_0$。\par \textbf{求导验算} 将原函数求导通分到$4V^2(x+1)$，分子化为4。\par \textbf{疑难点与方法比较} 本法无需一次设全部重因子的系数；只有模逆和低次部分分式。\end{example}

\subsection{有理积分的留数式系数法}\label{ic:indefinite_algebra_complete:residue}\begin{method}{有理积分的留数式系数法}\textbf{识别信号} 分母根已知，或只想先分离某几个极点。\par \textbf{操作步骤} 简单根$a$的系数为$\frac{N(a)}{D'(a)}$；$k$重根展开$\frac{(x-a)^kN}{D}$的前$k$项Taylor系数；共轭根成对还原实式。\par \textbf{疑难点与条件} 这里只计算局部部分分式系数，不涉及围道；和定积分的留数定理用途应区分。\end{method}

\begin{example}{竞赛风格｜重极点的局部展开}{难度：$\star\star\star$}\label{ic:ex:indefinite_algebra_complete:localcoeff}\textbf{题目} 分解并积分$\frac{1}{(x-1)^2(x+2)}$。\par \textbf{思路} 先在$x=1$展开剩余因子$\frac{1}{x+2}$。\par \textbf{详细解答} 写$x=1+h$，$\frac{1}{x+2}=\frac{1}{3}-\frac{h}{9}+O(h^2)$，故重极点贡献$\frac{1}{3(x-1)^2}-\frac{1}{9(x-1)}$。$x=-2$的简单项系数$\frac{1}{9}$，完整分解再加$\frac{1}{9(x+2)}$。原函数$-\frac{1}{3(x-1)}-\frac{\ln|x-1|}{9}+\frac{\ln|x+2|}{9}+C$。\par \textbf{求导验算} 求导后按上述完整分解通分为原式。\par \textbf{疑难点与方法比较} 局部展开避免大型待定系数方程，仍须检验全局真分式无余项。\end{example}

\subsection{平方分母的商导数降阶}\label{ic:indefinite_algebra_complete:quotient}\begin{method}{平方分母的商导数降阶}\textbf{识别信号} $\int \frac{f}{g^2}$且$\frac{f}{g'}$比原式简单。\par \textbf{操作步骤} 用$\dd(-\frac{1}{g})=\frac{g'\dd x}{g^2}$分部，得到$\int \frac{f}{g^2}=-\frac{f}{gg'}+\int\frac{f'g'-fg''}{g(g')^2}\dd x$；若$f=u'v-uv'$且$g=v$，直接识别商导数。\par \textbf{疑难点与条件} 须在$g,g'\ne0$的区间使用该变形；它是分部恒等式，不是极限的L’Hôpital法则。\end{method}

\begin{example}{竞赛风格｜藏在平方分母中的商}{难度：$\star\star\star$}\label{ic:ex:indefinite_algebra_complete:wronskian}\textbf{题目} 求$\int\frac{[x\cos x-\sin x+x^2\sin x]\dd x}{(x\cos x-\sin x)^2}$。\par \textbf{思路} 分子是$x$与分母的商导数交叉乘积。\par \textbf{详细解答} 设$g=x\cos x-\sin x$，则$g'=-x\sin x$，分子恰为$g-xg'$，因此原函数为$\frac{x}{g}+C=\frac{x}{x\cos x-\sin x}+C$。另一法把被积式拆成$\frac{1}{g}-\frac{xg'}{g^2}$，用分部积分$\int x\dd(\frac{1}{g})=\frac{x}{g}-\int\frac{\dd x}{g}$，两个剩余积分相消。\par \textbf{求导验算} $[\frac{x}{g}]'=\frac{g-xg'}{g^2}$，直接恢复原式。\par \textbf{疑难点与方法比较} 只在$g\ne0$的连通区间写同一个常数；商导数的交叉项符号必须按顺序核对。\end{example}

\subsection{反三角混合锁与有限维闭合}\label{ic:indefinite_algebra_complete:triglock}\begin{method}{反三角混合锁与有限维闭合}\textbf{识别信号} 指数的内层是$\arctan x$，分母含$(1+x^2)$的半整数幂。\par \textbf{操作步骤} 先令$t=\arctan x$，根式变为$\sec t$，再检查是否化指数三角积分；或以双元分部构造两个积分的线性关系。\par \textbf{疑难点与条件} 选择$t\in(-\frac{\pi}{2},\frac{\pi}{2})$保证$\sec t>0$；不要直接把平方根换成带任意符号的正割。\end{method}

\begin{example}{竞赛风格｜混合锁型的两解}{难度：$\star\star\star$}\label{ic:ex:indefinite_algebra_complete:atanlock}\textbf{题目} 求$\int \frac{x\e^{\arctan x}\dd x}{(1+x^2)^{\frac{3}{2}}}$。\par \textbf{思路} 换成角变量后是$\e^t\sin t$。\par \textbf{详细解答} $t=\arctan x$给$\int\e^t\sin t\dd t=\frac{\e^t(\sin t-\cos t)}{2}+C$，回代为$\frac{\e^{\arctan x}(x-1)}{2\sqrt{1+x^2}}+C$。另一法令$v=\sqrt{1+x^2}$，对$\frac{\e^{\arctan x}}{v}$分部；再对$\frac{x}{v}$用$(\frac{x}{v})'=\frac{1}{v^3}$，两式联立可解出同结果。\par \textbf{求导验算} 导数分子为$(x-1)+(1+x)$，除以$2v^3$后为$\frac{x}{v^3}$并乘$\e^{\arctan x}$。\par \textbf{疑难点与方法比较} 混合锁的难点是先看见角变量，而不是展开复杂高阶导数。\end{example}

\subsection{指数锁与归一化变量}\label{ic:indefinite_algebra_complete:explock}\begin{method}{指数锁与归一化变量}\textbf{识别信号} 根式中有$\e^x-x^2$，或分母为指数与多项式之和。\par \textbf{操作步骤} 先提指数的正因子，再检查$x\e^{-\frac{x}{2}}$的微分；对$g=\e^x+P$使用$g'-g=P'-P$将多项式与指数拆开。\par \textbf{疑难点与条件} 先写实定义域；指数提根式没有符号问题，但剩余根式仍须非负。\end{method}

\begin{example}{竞赛风格｜指数根式锁}{难度：$\star\star\star$}\label{ic:ex:indefinite_algebra_complete:exprootlock}\textbf{题目} 求$\int\frac{(x-2)\dd x}{\sqrt{\e^x-x^2}}$，在$\e^x>x^2$的区间。\par \textbf{思路} 提$\e^{\frac{x}{2}}$后，$t=x\e^{-\frac{x}{2}}$的导数完整出现。\par \textbf{详细解答} $\sqrt{\e^x-x^2}=\e^{\frac{x}{2}}\sqrt{1-t^2}$，$\dd t=-\frac12(x-2)\e^{-\frac{x}{2}}\dd x$，故答案$-2\arcsin(x\e^{-\frac{x}{2}})+C$。\par \textbf{求导验算} 链式求导给$\frac{(x-2)\e^{-\frac{x}{2}}}{\sqrt{1-x^2\e^{-x}}}$，化回原式。\par \textbf{疑难点与方法比较} 看似含指数和代数根式的混合题，可以由一个归一化变量解锁。\end{example}

\begin{example}{竞赛风格｜指数平方分母的闭合}{难度：$\star\star\star$}\label{ic:ex:indefinite_algebra_complete:expsquarelock}\textbf{题目} 求$\int \frac{x(x+1)\dd x}{(\e^x+x+1)^2}$。\par \textbf{思路} 记$g=\e^x+x+1$，有$g-g'=x$。\par \textbf{详细解答} $[\frac{x+1}{g}]'=-\frac{x\e^x}{g^2}$，而$\frac{x(x+1)}{g^2}=\frac{x}{g}-\frac{x\e^x}{g^2}$。又$\frac{x}{g}=1-\frac{g'}{g}$，故原函数$x-\ln|g|+\frac{x+1}{g}+C$。\par \textbf{求导验算} 导数$1-\frac{g'}{g}-\frac{x\e^x}{g^2}=\frac{x(g-\e^x)}{g^2}=\frac{x(x+1)}{g^2}$。\par \textbf{疑难点与方法比较} 先利用$g-g'$消掉指数，比对平方分母硬设全部系数更有针对性。\end{example}

\subsection{等分锁与对数微分配对}\label{ic:indefinite_algebra_complete:equalparts}\begin{method}{等分锁与对数微分配对}\textbf{识别信号} 分子是两个量的相对增长率之和，如$\frac{f'}{f}+\frac{g'}{g}$。\par \textbf{操作步骤} 先设$t=fg$或$\frac{f}{g}$，把分子看成$\dd\ln|t|$；再将分母中的同一$t$替换。\par \textbf{疑难点与条件} 相对微分只在非零区间成立；最终答案仍需检查原式定义域。\end{method}

\begin{example}{竞赛风格｜等分性转化为一元有理积分}{难度：$\star\star\star$}\label{ic:ex:indefinite_algebra_complete:equalpartlock}\textbf{题目} 求$\int\frac{(1+\cot x)\dd x}{1+\e^x\sin x}$。\par \textbf{思路} $t=\e^x\sin x$满足$\frac{\dd t}{t}=(1+\cot x)\dd x$。\par \textbf{详细解答} 积分$\int\frac{\dd t}{t(1+t)}=\ln|t|-\ln|1+t|+C=x+\ln|\sin x|-\ln|1+\e^x\sin x|+C$。\par \textbf{求导验算} 对两对数求导，公因子$1+\cot x$乘$1-\frac{t}{1+t}$，得到原式。\par \textbf{疑难点与方法比较} 等分性是对数链式法则；不需要把它当作没有条件的新积分规则。\end{example}

\subsection{代数原函数待定与伪椭圆判别}\label{ic:indefinite_algebra_complete:exactalg}\begin{method}{代数原函数待定与伪椭圆判别}\textbf{识别信号} 被积式$\frac{R(x)}{\sqrt{D(x)}}$，但有可能是$A(x)\sqrt D$或对数的导数。\par \textbf{操作步骤} 先设有理函数$A$，解$DA'+\frac{D'A}{2}=R$；对极点和无穷处比较次数确定候选形式；剩余项再试Pell型对数或真正椭圆标准形。\par \textbf{疑难点与条件} 某个候选失败不能证明没有初等原函数；伪椭圆指可初等化的具体微分，不是只看根式次数。\end{method}

\begin{example}{竞赛风格｜代数恰当微分的完整配系数}{难度：$\star\star\star$}\label{ic:ex:indefinite_algebra_complete:algebraansatz}\textbf{题目} 求$\int\frac{(1-x^4)\dd x}{(x^4+x^2+1)^{\frac{3}{2}}}$。\par \textbf{思路} 分母三次半幂提示原函数为低次多项式除以$\sqrt D$。\par \textbf{详细解答} 设$F=\frac{ax+b}{\sqrt D}$，则$F'=\frac{aD-\frac{(ax+b)D'}{2}}{D^{\frac{3}{2}}}$。比较奇偶项得$b=0$，再比较常数与$x^4$得$a=1$，故$F=\frac{x}{\sqrt{x^4+x^2+1}}+C$。\par \textbf{求导验算} 分子$D-\frac{xD'}{2}=1-x^4$。\par \textbf{疑难点与方法比较} 先用奇偶性和次数把待定式限制到最小，再解系数；不随意猜高次函数。\end{example}

\subsection{四次根式的多项式连分式算法}\label{ic:indefinite_algebra_complete:continuedfraction}\begin{method}{四次根式的多项式连分式算法}\textbf{识别信号} $D$为首项系数1的四次多项式；在无穷远作形式Laurent展开寻找对数型原函数。\par \textbf{操作步骤} 令$P_0=0,Q_0=1$，$a_k$取$\frac{\sqrt D+P_k}{Q_k}$的多项式部分。递推$P_{k+1}=a_kQ_k-P_k$、$Q_{k+1}=\frac{D-P_{k+1}^2}{Q_k}$。收敛分式$\frac{p_k}{q_k}$按$p_k=a_kp_{k-1}+p_{k-2}$、$q_k=a_kq_{k-1}+q_{k-2}$计算。每步检验$p_k^2-Dq_k^2$是否为非零常数。\par \textbf{疑难点与条件} 初值$p_{-2}=0,p_{-1}=1,q_{-2}=1,q_{-1}=0$。有限步未找到恒等式时，不能凭此宣布不可积；形式展开不依赖原函数的实定义区间。\end{method}

\begin{theorem}{Pell恒等式的最终证书}设$P,Q,D$为实多项式，在所用区间有$D>0,Q\ne0$。若$P^2-DQ^2=c\ne0$，则$[\ln|P+Q\sqrt D|]'=\frac{P'}{Q\sqrt D}$。若进一步$\frac{P'}{Q}=\alpha x+\beta$，就得出该线性分子与四次根式的原函数。\end{theorem}\begin{proof}微分恒等式得$PP'=DQQ'+\frac{D'Q^2}{2}$。对数导数的分子$P'+Q'\sqrt D+\frac{QD'}{2\sqrt D}$等于$\frac{P'(P+Q\sqrt D)}{Q\sqrt D}$，约分得到结论。代数证书必须最后独立展开核对。\end{proof}

\begin{example}{竞赛拓展｜连分式积分的五步计算}{难度：$\star\star\star$}\label{ic:ex:indefinite_algebra_complete:continuedfull}\textbf{题目} 求$\int \frac{x\dd x}{\sqrt{(x+1)^4-12x^2}}$，在根式为正的区间。\par \textbf{思路} 用多项式连分式寻找常数Pell余项；本例可在第4号收敛分式找到。\par \textbf{详细解答} 记$D=(x+1)^4-12x^2$。按上述普通连分式递推，得到下表。\par \noindent\begin{tabular}{ccl}\toprule $k$&$a_k$&$p_k^2-Dq_k^2$\\\midrule 0&$x^2+2x-5$&$-24(x-1)$\\1&$\frac{x+3}{12}$&$-\frac{x}{3}$\\2&$-6(x+2)$&$-36x$\\3&$\frac{x+2}{18}$&$-\frac{2(x-1)}{9}$\\4&$-9(x+3)$&$108$\\\bottomrule\end{tabular}\par 将最后分式同时乘$-4$，得\begin{align*}P&=x^6+12x^5+45x^4+44x^3-33x^2+43,\\Q&=x^4+10x^3+30x^2+22x-11.\end{align*}独立展开验证$P^2-DQ^2=1728$、$P'=6xQ$，所以所求为$\frac16\ln|P+Q\sqrt D|+C$。\par \textbf{求导验算} Pell对数微分给$\frac{P'}{6Q\sqrt D}=\frac{x}{\sqrt D}$。若$Q$在某点为0，最后对数式可按原被积函数的正则性连续延拓，而中间商式仅在$Q\ne0$时使用。\par \textbf{疑难点与方法比较} 换元后的根式必须重新通分核验。本例以显示的$D$和两条多项式恒等式独立验证结果，所有幂次均由展开确定。\end{example}

\subsection{微分方程引领积分：Lagrange恒等式}\label{ic:indefinite_algebra_complete:odeidentity}\begin{method}{微分方程引领积分：Lagrange恒等式}\textbf{识别信号} 已知$y''+py'+qy=0$，被积函数与$y$及其导数有关。\par \textbf{操作步骤} 选简单$h$，计算$h''+ph'+qh$；将目标凑成这个组合乘$y\e^{\int p}$，再写出边界组合$h'y-hy'$。\par \textbf{疑难点与条件} $h$须有二阶导数，系数和解在同一区间正则；不是任意选$h$都能匹配目标。\end{method}

\begin{theorem}{二阶方程积分公式}若$y''+py'+qy=0$，则$\int \e^{\int p}(h''+ph'+qh)y\dd x=\e^{\int p}(h'y-hy')+C$。\end{theorem}\begin{proof}对右端求导，得到$\e^{\int p}[(h''+ph')y-h(y''+py')]$，用原方程$y''+py'=-qy$即为左端被积式。\end{proof}

\begin{example}{竞赛风格｜ODE恒等式与分部互证}{难度：$\star\star\star$}\label{ic:ex:indefinite_algebra_complete:ode_sine}\textbf{题目} 求$\int(x^2+2)\sin x\dd x$。\par \textbf{思路} $\sin x$满足$y''+y=0$，取$h=x^2$。\par \textbf{详细解答} $h''+h=x^2+2$，故原函数$h'y-hy'=2x\sin x-x^2\cos x+C$。两次分部也给$\int x^2\sin x=-x^2\cos x+2x\sin x+2\cos x$，加$2\int\sin x$后余项消去。\par \textbf{求导验算} 导数中的$2x\cos x$互相抵消，余下$(x^2+2)\sin x$。\par \textbf{疑难点与方法比较} ODE恒等式可先识别整组外力，避免逐项积分后再相消。\end{example}
